\section{The Cognitive and Computational Basis of Appraisal}
\label{p11:sec:appraisal}

\noindent
Before evaluation is an explicit measurement, it is a continuous and largely unconscious inference. The parties to a shared life do not, for the most part, assess their relation by deliberate procedure; they are appraising it constantly, beneath reflection, registering in the texture of an ordinary day whether things are well or ill between them, whether the relation is tending toward warmth or toward distance, whether the shared life is rising or declining. This section establishes that appraisal is, at its basis, an inference upon a partially observable state, conducted in a value that is not symbolised; that the explicit, instrumented evaluation of \xref{p11:sec:evaluation} is a structured proxy laid over this basal inference; and that the framework which best renders the basal process is the framework of belief updating, which is evidentialism reaching below the deliberate weighing of studies to the continuous revision of a belief about an unobservable state.

\subsection{Appraisal as a biological process}
\label{p11:sec:ap-bio}

The appraisal of a situation is, in the first instance, a biological event and not a deliberation. The tradition of appraisal theory in the psychology of emotion holds that emotion arises from an evaluation of a situation with respect to a being's concerns, and that this evaluation runs through rapid, multilevel, largely automatic checks before any of it becomes available to reflection \parencite{lazarus1991,scherer2001}. That emotion and cognition are not separable faculties but aspects of one evaluative process is argued by \textcite{pessoa2008}, and the constructed and inferential character of emotional appraisal by \textcite{barrett2017}. The felt quality of a moment in a shared life, the ease or the disquiet that arrives unbidden after an exchange, is the registration of such an appraisal, an evaluation already accomplished beneath awareness and delivered to consciousness as a feeling. The neuroscience of this process locates it in the rapid integration, by prefrontal and associated circuits, of present perception with prior expectation, and the issuing of an evaluation of the current state on which decision is then based. Crucially, the process operates on quantities the subject does not and cannot fully articulate: the appraisal proceeds in a value that is felt and acted upon before, and often without ever, being symbolised. The basal appraisal is conducted in exactly the unsymbolised value the prior papers held to be the real one.

A component of this process deserves emphasis, because it grounds the connection to the matter of the field. Interoception, the sensing of the body's internal condition, is increasingly understood as itself a process of inference, in which the felt bodily state is the brain's best estimate of the causes of its internal signals, and emotion a reading of this estimate in context \parencite{seth2013}. The felt sense of how a relation stands, the bodily ease or unease that is the first and most constant appraisal of a shared life, is on this account an interoceptive inference, and its data are in large part the sensory and habitual fabric the prior paper identified as the matter of the field: the meals, the rooms, the rhythms, the contact, registered in the body before they are formulated in thought. The appraisal of a relation is, at its basis, the body's continuous inference upon the texture of the shared life.

\subsection{Appraisal as inference upon a partially observable state}
\label{p11:sec:ap-pomdp}

The formal structure of this basal appraisal is rendered by the framework of decision under partial observability. An agent that cannot observe the true state of its world directly, but only through partial and noisy observations, must maintain a belief, a probability distribution over the possible states, and update this belief as observations arrive; it then chooses its actions on the belief rather than on the unobtainable true state. This is the structure formalised as the partially observable Markov decision process, in which a hidden state evolves, the agent receives observations that depend on it, the agent maintains and updates a belief over the hidden state, and the agent selects actions to secure value over time \parencite{kaelbling1998}. The general framework of learning to act for value under uncertainty is set out in \textcite{sutton_barto_rl}.

The structure maps onto the appraisal of a relation with a precision that is instructive, and the mapping carries the commitments of the prior papers into the vocabulary of decision.

\begin{claim}[The appraisal of a relation as inference upon a hidden state]
The basal appraisal of a relation has the structure of inference upon a partially observable state. The true state of the relation is hidden, not observable directly, in keeping with the prior result that the real value of a relation is transversal across the registers and not exhausted by any observation. What is observed is the everyday sensory and habitual fabric of the shared life, the matter of the field, which depends on the hidden state without revealing it fully. Each party maintains a belief, not a determinate verdict but a distribution, over the state of the relation, in keeping with the prior account of relational value as superposed rather than definite. The belief is updated as the fabric is lived, and action, which is practice, is taken upon the belief. The appraisal is thus an inference, conducted in an unsymbolised value, upon a state that no observation discloses in full.
\end{claim}

\noindent
The mapping has a consequence that connects this section to the formal dilemma of the prior paper. A standard process of this kind operates with a fixed value function, a fixed specification of what the agent is to secure. The appraisal of a relation does not. What the parties are to secure, the good of the relation, is itself reconceived across the lifecycle, at the moment of re-cognition (\xref{p11:sec:recognition}); the value function is rewritten by the process it governs. This is the reflexive reconfiguration of the prior papers, and it reappears here as the non-standard feature that the value function of the relational decision process is not fixed but is revised by the cycle, which is the decision-theoretic form of the thesis that there is no fixed meta-value-grammar.

\subsection{Self-maintenance and the imperative to keep inferring}
\label{p11:sec:ap-active}

A development of the inferential account binds appraisal to the survival of the relation. On the free-energy account of the functioning of adaptive systems, an organism maintains itself, holds itself away from the disordered states that would dissolve it, precisely by continually minimising the discrepancy between its expectations and its sensory states, which is to say by continually inferring and acting to make its inferences good \parencite{friston2010}. The process-theoretic development of this principle, in which perception and action jointly minimise expected surprise, is given by \textcite{friston2017}. Perception, action, and the maintenance of the organism's own existence are, on this account, one process: to cease to infer and to act upon the inference is to cease to maintain the boundary that constitutes the system, and so to dissolve.

Transposed to the relation, this account supplies the deep form of a thesis the paper will develop in dynamical terms (\xref{p11:sec:regimes}). The continual appraisal of a relation, the unbroken updating of each party's belief about its state and the action taken upon it, is not an optional accompaniment to the relation but the activity by which the relation maintains itself. A relation in which the parties cease to appraise, cease to update their sense of one another, cease to act upon that sense, is a relation that has ceased to perform the activity by which it holds itself in being, and it dissolves toward the terminal rest of \xref{p11:sec:dyn-fixed}. The imperative to keep inferring, on this account, is the imperative to keep the relation in existence; the cessation of appraisal is the beginning of extinction.

\subsection{The two registers of evaluation}
\label{p11:sec:ap-two}

The foregoing establishes a distinction that governs the paper's treatment of evaluation. There are two registers in which a relation is evaluated. The first is the basal, continuous, largely unconscious inference described in this section: an appraisal conducted in unsymbolised value, upon a hidden state, through the lived fabric of the shared life, and running through every moment of the cycle without occupying any one of them. This is evaluation in its primary and pervasive form, and it is conducted in the real value the prior papers defended. The second is the explicit, instrumented, structured evaluation that the positivist framework supplies and that \xref{p11:sec:evaluation} develops: a deliberate measurement, conducted at particular moments, in which some aspect of the relation is brought to a structured and assessable form.

The relation between the two registers is the relation the operational concession named (\xref{p11:sec:bg-concession}). The explicit evaluation is a structured proxy laid over the basal inference: it takes the unsymbolised belief that the parties continuously maintain and renders some structured aspect of it in a form that can be measured, compared, and accumulated. The proxy is operationally valuable, since the basal inference, precisely because it is unsymbolised, is private, incommunicable, and unaccountable, while the structured measure is public and checkable. And the proxy is incomplete, since the structuring discards exactly the unsymbolised remainder in which the real appraisal is conducted. To evaluate a relation explicitly is to symbolise, for operational use, a belief that is in its nature unsymbolised, and the gain in publicity is paid for in a loss of the real. The paper's treatment of explicit evaluation in \xref{p11:sec:evaluation} proceeds under this understanding throughout: the measures it discusses are structured proxies for a basal appraisal they do not reach, and evidentialism, which is the epistemology of the basal belief updating itself, is the framework within which both the basal inference and its structured proxy are to be understood.
